\documentclass[11pt]{article}
\usepackage[margin=0.68in]{geometry}
\usepackage{fontspec}
\usepackage{setspace}
\usepackage{hyperref}
\hypersetup{colorlinks=true, urlcolor=blue, linkcolor=black}
\pagestyle{empty}
\setlength{\parskip}{0.38em}
\setlength{\parindent}{0pt}
\setstretch{0.91}
\setmainfont{Times New Roman}

\begin{document}

\begin{flushleft}
Ke Cheng HSU (Clio) \\
Department of Economics, The Hong Kong University of Science and Technology \\
Clear Water Bay, Kowloon, Hong Kong $\cdot$ +852 6840 9753 \\
\href{mailto:khsu@connect.ust.hk}{khsu@connect.ust.hk} $\cdot$ \href{https://cliohsu1997.github.io/}{cliohsu1997.github.io}
\end{flushleft}

28 November 2026

\vspace{0.02em}
Faculty Search Committee \\
Shenzhou Business School \\
Eastern Institute of Technology, Ningbo \\
Ningbo, China

\vspace{0.02em}
Dear Members of the Search Committee,

I am writing to apply for an Assistant Professor position in Business Analytics at Shenzhou Business School. I learned of the opening from the AIS Academic Careers posting. I am a Ph.D.\ candidate in Economics at the Hong Kong University of Science and Technology. I expect to graduate in June 2027 and can start in July 2027. My fields are Empirical Industrial Organization and Digital Economics. I am intrigued by the School's emphasis on cooperating with industry and on helping firms solve real problems. That is already how my research is built: each paper starts from a decision a firm has to make, and the answer is estimated on the firm's own data.

My job market paper, ``Designing Return Policy in E-Commerce: A Trade-off between Purchase and Return Efficiency,'' studies whether a return policy encourages efficient returns and purchases. A stricter policy curbs inefficient returns after delivery, but the lower return rate makes a higher product price more attractive for the seller and can exclude efficient purchases. Using transaction data from a Taobao apparel retailer, I estimate a structural model of purchase, return, and whether to buy return-shipping insurance under monopoly pricing across six categories. Automatic coverage of return shipping performs better when the margin, net scrap value, and uncertainty are high, because it expands efficient purchases. Optional coverage performs better when those attributes are low, because it curbs excessive returns. A menu that pairs a discounted price with a medium restocking fee balances the two margins and performs better than both insurance policies for most categories.

A second working paper, ``Product Design with Heterogeneous Search Costs'' (with Pak Hung Au), studies how firms choose product design when consumers differ in search costs. Search costs shape which products consumers search. Consumers with low search costs seek niche products, and consumers with high search costs seek standardized products. A fall in search costs can raise prices when this shift in which product consumers search outweighs the effect that encourages more search. I also have two ongoing projects. With Kohei Kawaguchi and Kei Kawai, I use Japanese bookstore data to study how resale price maintenance encourages retailers to hold inventory in low-demand states. However, unsold copies cannot be cleared by cutting the price, so they are returned and discarded at an environmental cost. With Ce Bian and Liang Zhao, I study pass-through when products can be returned, and show that a tax can lower the equilibrium price and raise welfare by reducing the return rate.

At HKUST I was a teaching assistant from 2020 to 2025 for undergraduate macroeconomics, the economics of uncertainty and information, and Ph.D.\ microeconomic theory, with an average evaluation of 4.72/5. In 2026 I designed and lectured the mathematics camp for the MSc in Economics, receiving 4.53/5. Eastern Institute of Technology's published undergraduate curriculum in Big Data Management and Application is organized around data analysis and business intelligence. The courses in that list I can teach are Data Statistics and Econometric Analysis, Data Mining and Business Applications, and Data Intelligence and Decision Analysis, at undergraduate level and in the corresponding postgraduate courses, with little additional preparation. A data-mining course of this kind covers regression, classification, clustering, and recommendation systems, which I implement in R, Python, and C++. Data Statistics and Econometric Analysis is the applied-econometrics course in that curriculum. After probability and statistics, it covers regression, the potential-outcome model, difference-in-differences, instrumental variables, panel data, and discrete choice. That is the econometrics sequence in my doctoral training. I am prepared to teach these courses in English.

My referees' contact information is in my CV. Crystal, the secretary of the Department of Economics at HKUST, will send letters of reference from Pak Hung Au, Kohei Kawaguchi, and Kei Kawai to \href{mailto:SHENZHOU@eitech.edu.cn}{SHENZHOU@eitech.edu.cn}, with a copy to \href{mailto:faculty-hire@eitech.edu.cn}{faculty-hire@eitech.edu.cn}. You can reach me at \href{mailto:khsu@connect.ust.hk}{khsu@connect.ust.hk} or +852 6840 9753.

\vspace{0.05em}
Sincerely,

\vspace{0.05em}
Ke Cheng HSU \\
Ph.D.\ Candidate, Department of Economics, HKUST

\end{document}
